\documentclass[11pt]{article}

\usepackage[
	top=20mm,
	bottom=20mm,
	left=20mm,
	right=20mm,
	marginparwidth=1cm,
	marginparsep=1mm
]{geometry}

%\setlength{\parskip}{0.5\baselineskip}%
%\setlength{\parindent}{0pt}%

\usepackage{fancyhdr}
\renewcommand{\headrulewidth}{.4mm} % header line width
\pagestyle{fancy}
\fancyhf{}
\fancyhfoffset[L]{1cm} % left extra length
\fancyhfoffset[R]{1cm} % right extra length
\lhead{MAT137Y Instructor Guide}
\rhead{Course Goals}
\rfoot{}
\cfoot{\thepage}

%%%%%%%%%%%%%%%%%%%%%%%%%%%%%%%%%%
%%%%%%%		PACKAGES
%%%%%%%%%%%%%%%%%%%%%%%%%%%%%%%%%%
\usepackage{setspace} % controlling line spacing
\setlength{\parindent}{0pt} % paragraphs are not indented
\usepackage{amssymb}
\usepackage{graphicx}
\usepackage{enumitem}
\usepackage{amsfonts}
\usepackage{amssymb}
\usepackage{ifthen}
\usepackage{multicol}

\usepackage{tikz}
\usetikzlibrary{shapes, backgrounds}

\usepackage[english]{babel}

\setlength{\parindent}{0cm}
%\setlength{\parskip}{1em}

\newcommand{\DS}[1]{${\displaystyle #1}$}
\newcommand{\vv}{\vspace{.5cm}}
\newcommand{\n}{\newpage}

\newcommand{\R}{\mathbb{R}}
\newcommand{\Q}{\mathbb{Q}}
\newcommand{\Z}{\mathbb{Z}}
\newcommand{\N}{\mathbb{N}}
\newcommand{\floor}[1]{\lfloor #1 \rfloor}

\definecolor{gold}{rgb}{212,175,55}
\definecolor{mygreen}{RGB}{66,165,81}
\newcommand{\azul}[1]{{\color{blue} #1}}
\newcommand{\rojo}[1]{{\color{red} #1}}
\newcommand{\verde}[1]{{\color{mygreen} #1}}

\newcommand{\set}[2]{ \left\{ #1 \; : \; #2 \right\} }
\newcommand{\e}{\varepsilon}

\definecolor{137cp1}{RGB}{13, 33, 161}
\definecolor{137cp2}{RGB}{51, 161, 253}
\definecolor{137cp3}{RGB}{255, 67, 101}%{255, 84, 0}
\definecolor{137cp4}{RGB}{232, 144, 5}%{255, 190, 11}%{83, 221, 108}%{38, 196, 133}

\usepackage{hyperref}
\hypersetup{colorlinks}
\hypersetup{urlcolor=137cp3, linkcolor=137cp1}

\usepackage{sectsty}
\sectionfont{\centering \LARGE\color{137cp3}} % sets colour of chapters
\subsectionfont{\Large \color{137cp2}}
\subsubsectionfont{\large \color{137cp4}}
\paragraphfont{\color{137cp1}}

\usepackage{tikzsymbols}
\usepackage[final]{pdfpages} %insert .pdf file

\setcounter{secnumdepth}{0}
\setcounter{tocdepth}{2}

\usepackage{amsthm, thmtools}
\usepackage{mdframed}

%% kill warnings for overfull hboxes
%\newcommand{\ignoreoverfullhboxes}{\setlength{\hfuzz}{\maxdimen}}
%\AtBeginEnvironment{mdframed}{\ignoreoverfullhboxes}

%==========================================
%: theorem styles
%==========================================

\declaretheoremstyle[
spaceabove=-6mm,
spacebelow=-2cm,
headfont=
\color{137cp1}
\bfseries,
notefont=\bfseries\mathversion{bold},
notebraces={(}{)},
%bodyfont=\itshape,
postheadspace=2mm,
headpunct={.}
]{myexample}

\declaretheoremstyle[
spaceabove=-6mm,
spacebelow=-2cm,
headfont=
\color{137cp3}
\normalfont,
bodyfont=\normalfont,
postheadspace=2mm,
headpunct={.}
]{mywarning}

%==========================================
%: theorem environments
%==========================================
% ??? create theorem synopsis - list of only important theorems/prop/lem - environments with option to add to this list

\definecolor{Lavender}{rgb}{0.95,0.90,1.00}
\definecolor{darkviolet}{rgb}{0.35,0.00,0.70}

\newcommand{\mypartscolour}{Lavender!50}
\newcommand{\mylemmacolour}{darkviolet!15}
\newcommand{\mydefinitioncolour}{red!50}
\newcommand{\myremarkcolour}{yellow!25}

%: 	EXAMPLE

\declaretheorem
[style=myexample,
name=Exercise,
refname={Exercise},
Refname={Exercise},
]
{thmx}

\DeclareDocumentEnvironment{exercise}{O{ } g} % optional arguments: title, label
{\begin{mdframed} [backgroundcolor=\mypartscolour, skipabove=0.5\baselineskip, innertopmargin=0.5\baselineskip, skipbelow=0.5\baselineskip, innerbottommargin=0.1\baselineskip, roundcorner=12pt leftmargin=0.25cm, rightmargin=-0.25cm, innerleftmargin=0.3cm, innerrightmargin=0.25cm, linewidth=4pt, linecolor=137cp1, hidealllines=true, leftline=true, nobreak=false, roundcorner=50pt ] \begin{thmx}[#1] \IfNoValueTF{#2}{}{\label{#2}\hypertarget{#2}{}}}
{\end{thmx} \end{mdframed}}

%: 	WARNING
\declaretheorem
[style=mywarning,
name=Warning,
numbered=no]
{propx}

\DeclareDocumentEnvironment{warning}{O{ } g} % optional arguments: title, label
{\reversemarginpar\marginpar{\hspace{10cm} \includegraphics[height=18pt]{alert.png} } \vspace{-3mm} \begin{mdframed} [backgroundcolor=red!10, skipabove=0.5\baselineskip, innertopmargin=0.5\baselineskip, skipbelow=1\baselineskip, innerbottommargin=0.5\baselineskip, leftmargin=-0.25cm, rightmargin=-0.25cm, innerleftmargin=0.25cm, innerrightmargin=0.25cm, linewidth=3pt, linecolor=137cp3, hidealllines=true, leftline=true, nobreak=false] \begin{propx}[#1]%
\IfNoValueTF{#2}{}{\label{#2}\hypertarget{#2}{}}}
{\end{propx} \end{mdframed} }

%%%%%%%%%%%%%%%%%%%%%%%%%%%%%%%%%%%%%%%%%%%%%%%%%%%%%%%%
%%%%%%%%%%%%%%%%%%%%%%%%%%%%%%%%%%%%%%%%%%%%%%%%%%%%%%%%
%%%%%%%%%%%%%%%%%%%%%%%%%%%%%%%%%%%%%%%%%%%%%%%%%%%%%%%%

\usepackage{amsthm, thmtools}
\usepackage{mdframed}

% kill warnings for overfull hboxes
\newcommand{\ignoreoverfullhboxes}{\setlength{\hfuzz}{\maxdimen}}
\AtBeginEnvironment{mdframed}{\ignoreoverfullhboxes}

\renewcommand{\baselinestretch}{1.2} %increase line spacing

\renewcommand{\labelitemi}{$\textcolor{137cp1}{\bullet}$}

%%%%%%%%%%%%%%%%%%%%%%%%%%%%%%%%%%%%%

\begin{document}
	\thispagestyle{plain}
	\begin{center}
		{\bf {\LARGE \textcolor{137cp3}{MAT137Y: Course Goals} }

		\medskip {\Large \textcolor{137cp1}{Alfonso Gracia-Saz} }}
	\end{center}


	This guide will help you understand better who our students are as well as the main goals of the course.

	The section \hyperref[who]{Who are our students?} gives a brief description of who MAT137Y students are: what is the difference between MAT137Y and other calculus courses offered at the University of Toronto, what to expect from our students, their background, course prerequisites, etc.

	The section \hyperref[GCO]{Global Course Objectives} describes the three major objectives in the course. These are very important: they are the main point of reference that we use when making pedagogical decisions.

	Finally, the \hyperref[Appendix]{Unit-by-Unit Course Objectives} section provides a very detailed list of the objectives we would like to achieve in each unit. You do not have to read all of these now. You may prefer to postpone them till you are planning your classes for each unit.



	\tableofcontents

	\newpage

	%%%%%%%%%%%%%%%%%%%%%%%
	\section{Who are our students?}

	\paragraph{MAT135/6 vs MAT137 vs MAT157.}
	\begin{itemize}
		\item Students in the most ``hardcore" math programs (Pure Math and three others)
			take \textcolor{137cp1}{MAT157Y}, an analysis course that sometimes begins
			with Dedekind cuts and it assumes they are already fluent in proofs and
			rigour (or that they will learn it by osmosis).

		\item Students in life science who need a general calculus course without proofs
			or formal definitions take \textcolor{137cp1}{MAT135H$+$MAT136H}. It focuses
			on modelling, applications, and communication.

		\item \textcolor{137cp1}{MAT137Y} is a hybrid in the middle.

			There are four math specialist programs that begin with MAT137Y. These are
			\begin{itemize}
				\item Mathematical applications in economics and finance.

				\item Mathematics and its applications in probability and statistics

				\item Mathematics and its applications in physical science

				\item Mathematics and its application in teaching
			\end{itemize}
			However, the majority of our students are not in a math specialist program.
			They are students in actuarial science, statistics, computer science,
			economics, or physics (or perhaps in a math major). Their programs may
			require MAT137Y or simply recommend it.
	\end{itemize}

	\paragraph{What to expect from our students.}
	Our student are smart, motivated, and were ``good at math in high school". Unfortunately,
	this does not mean much. Most of them have only taken math courses consisting in
	memorizing formulas and algorithms and plugging numbers into them. They may be
	comfortable with heavy computations, but they have never been asked to understand
	why a method works or to think creatively. They are capable of learning rigour,
	but they have never been exposed to it before.

	In MAT137Y we acknowledge that mathematical abstraction and proofs are difficult. We assume no background knowledge on proofs. We will teach all the related skills (see \hyperref[CO2]{Global Course Objective 2} below) slowly and from scratch. This is a big difference with an analysis course like MAT157Y.

	\paragraph{Prerequisites.}
	In terms of prerequisites, we assume that students have mastered precalculus
	topics (trigonometry, exponentials and logarithms, basic graphing, polynomial
	manipulations...) This is not always true, but students are responsible to do any
	catch up they need to do by themselves, and we provide
	\href{http://uoft.me/precalc}{http://uoft.me/precalc} for this purpose. There is
	one exception that we include and teach in MAT137: the proper definitions of
	function and inverse function (see \hyperref[unit4]{Unit 4} below).

	\paragraph{Goal.}
	Having talked to professors in the non-math departments who send their students
	to MAT137Y, their priority is that students accomplish
	\hyperref[CO3]{Global Course Objective 3} below.

	Finally, a good way to summarize the objectives of MAT137 is ``to get students ready to take MAT237Y". While nominally a multivariable calculus course, MAT237Y is even more a hybrid calculus/analysis. It studies more advanced concepts (such as compactness, uniform continuity, or completeness) and it requires students tor read, understand, and write much more complex proofs. 
	%%%%%%%%%%%%%%%%%%%%%%%
	\newpage

	\section{Global Course Objectives}


	There are three major objectives in the course:

	%%%
	\subsection{1. Calculus concepts}


	We want students to become fluent in the main calculus concepts (limits, derivatives,
	integrals, sequences, series, and power series). This includes:
	\begin{itemize}
		\item State their definitions precisely.

		\item Recognize and construct examples and non-examples.

		\item Perform simple computations.

		\item Remember the main theorems that involve these concepts, state them
			precisely, and use them. This includes at least the Intermediate Value
			Theorem, the Extreme Value Theorem, limit laws, differentiation rules, the
			Mean Value Theorem, L'H\^{o}pital's Rule, the Fundamental Theorem of Calculus,
			and the Monotone Convergence Theorem for sequences. Notice that the formal
			proofs of some of these are too complex and not relevant to this course.

		\item Use these concepts in simple proofs.

		\item Determine whether a claim involving them is true or false.
	\end{itemize}

	%%%
	\subsection{2. Mathematical rigour}

	We want to introduce students to mathematical rigour and logic. We want them to
	read and understand mathematical statements and precise definitions, as well
	as read, critique, and write rigorous proofs.

	There are a lot of small skills that are part of this objective. Since they are all second nature to us, it is easy to forget how foreign and difficult they are to most of our students. In order to acquire these skills, students need to explicitly practice all of them: to do them themselves, and to receive feedback. They do not learn any of this ``by osmosis" (by watching us demonstrate them). In particular ``watching a proof of theorem X" is not a learning objective.
	The list of skills related to this big objective include:
	\begin{itemize}
		\item Understand the difference between a definition and a theorem.

		\item Translate an intuitive concept or idea into a precise mathematical definition.

		\item Write the negation of a mathematical definition.

		\item Read a mathematical definition, and determine whether a particular
			object is or isn't an example.

		\item Compare mathematical statements and determine whether they are equivalent,
			whether one is strictly stronger, or whether they are not comparable

		\item Read a mathematical statement and determine whether it makes sense or it
			is nonsense. If it makes sense, determine whether it is true or false.

		\item Construct examples satisfying certain properties.

		\item Prove a claim to be false by constructing a counterexample (and know
			that this is the way to do it).

		\item Read and critique simple proofs; this includes finding the flaws in a
			bad proof and fixing them if possible.

		\item Write simple proofs. This could mean reconstruct a proof they have
			seen before, use the same ideas in a different context, or write a new
			proof. At this level, students only need to write \emph{simple} proofs (what
			we normally call ``one-liners").

		\item Use set notation, set-building notation, quantifiers, and conditionals
			in the context of all the above.

		\item Understand and justify why a proof needs to have a certain structure,
			why we need to write things in a certain way, and why we consider some things
			a valid proof and others not.

			\begin{quote}
				In other words, memorizing templates for how to write proofs is not
				enough. For example, as long as they are always asked questions with the
				same format, it is possible for a student to memorize how to write a ``perfect"
				proof by induction and having no idea of why proof by induction works. This
				does not meet our objectives.
			\end{quote}

		\item Use variables in proofs correctly. It is hard to overestimate how
			difficult this is for students. They need to practice it (and receive
			feedback on it) as much as possible. We need to be strict with it to
			properly relay its importance. Specifically, students need to know, and correctly
			use in proofs:
			\begin{quote}
				Variables in a definition or claim normally have to be quantified. A
				quantified variable is a ``dummy variable" and it does not carry any intrinsic
				meaning. Variables in a proof should not be quantified; they should be fixed.
				A variable cannot be used until it has been fixed. Variables need to be introduced
				in order, paying attention to what depends on what.
			\end{quote}
	\end{itemize}

	As examples, here are two of the most difficult claims that students should be
	able to prove.

	\begin{exercise}
		{} Let $a \in \mathbb{R}$. Let $f$ be a function defined at least on an
		interval centered at $a$, except possibly $a$. Prove that
		\[
			\text{ IF }\lim_{x \to a}f(x) = \infty, \quad \quad \text{ THEN }\lim_{x
			\to a}\frac{1}{f(x)}= 0.
		\]
		Write a proof directly from the definitions of limit, without using any of the
		limit laws.
	\end{exercise}

	\begin{exercise}
		{} Let $a <b.$ Let $f$ be a function differentiable on $(a,b)$. Prove that
		\[
			\text{IF for all }x \in (a,b), \, f'(x)>0, \quad \quad \text{ THEN }f \text{
			is increasing on }(a,b).
		\]
		\emph{Hint:} Use the Mean Value Theorem.
	\end{exercise}

	Proofs more difficult than these two exercises are probably inappropriate in a
	test in MAT137Y, although they might be suitable as a learning exercise with
	lots of scaffolding. If students really understand how proofs and rigour work using
	simple results in MAT137Y, then they can move on to other courses (MAT237Y)
	where they will learn to write longer and more complex proofs.

	As a warning, we cannot expect students to acquire all these skills in \hyperref[unit1]{Unit 1} (``Intro to logic and proofs"). \hyperref[unit1]{Unit 1} provides students with the basics so that they begin to understand what a proof is and they have a chance to start practicing. However, they will only master proof writing through consistent practice (and feedback) throughout the whole course. 

	%%%
	\subsection{3. Problem solving}

	We want students to be comfortable attacking new problems that they have not
	seen before, to make conjectures, to figure out by themselves how to adjust
	old methods to new situations, to figure out what new things they need to learn
	and to learn them, and to know when to be confident with their answers. We want
	students to build resilience, persistence, and creativity. We want them to
	become comfortable with struggle and productive failure.

	To accomplish this, whenever possible (in class and in assignments) we make students part of the discovery process and we guide them to come up with claims, algorithms, or proofs themselves. We emphasize building intuition and understanding why a method or algorithm or theorem works (which is not the same as having seen a rigorous proof for it, and often has nothing to do with it).

	Whenever possible we avoid asking students to memorize formulas or patterns and use them mindlessly. That hurts rather than helps in our objective.

	%%%%%%%%%%%%%%%%%%%%%%%%%%

	\vfill

	\begin{quotation}
		\noindent
		\emph{``I thoroughly enjoyed this course, as my first real introduction to math
		as not just a list of formulas, but indeed a exploration of abstract
		mathematical objects and their properties. It helped me discover a love of math
		I did not know that I had, and definitely has motivated me to explore a
		higher level of math than I planned to in my time at U of T. I did not consider
		myself particularly strong at math in high school as I am not very good at memorizing
		formulas (high school trigonometry was an unmitigated disaster), but your
		course empowered me as a student of math in a way that I haven't been before,
		as I came to appreciate the abstract logic underlying each and every concept,
		from basic limit definitions all the way to Taylor series applications."}

		\hfill --- From a student's unsolicited email at the end of the course
	\end{quotation}

	%%%%%%%%%%%%%%%%%%%%%%%
	\newpage
	\section{Unit-by-Unit Course Objectives}


	\begin{warning}
		The detailed list below is intended to help, but it is much less important
		than the global course objectives above. It is also much more dynamic; it
		can and will likely evolve.
	\end{warning}

	\setcounter{subsection}{0}

	%%%
	\subsection{1. Introduction to logic, sets, notation, definitions, and proofs}


	The goal of \hyperref[unit1]{Unit 1} is to provide students with the basics so
	that they begin to understand notation, definitions, and proofs and they have a
	chance to start practicing all the skills listed in \hyperref[CO2]{Global
	Course Objective 2}.

	%%%
	\subsection{2. Limits and continuity}


	When I write ``limit" it means all ``versions" of limits (finite, infinite, at
	a point, at infinity, side limits).

	\begin{itemize}
		\item Understand the concepts of limit intuitively, and estimate limits from
			graphs and tables of values.

		\item Write formal definitions of limits. Determine which variations of the
			definition are equivalent and otherwise determine what they mean.

		\item Prove that limits exists or do not exist from the definition.

		\item State and prove the limit laws and the Squeeze Theorem. Construct
			proofs of other similar results. Use them in computations and in other simple
			proofs.

		\item Define continuity, identify continuous functions, and the various types
			of discontinuities.

		\item Understand the theorems that are or are not true about composition of limits
			and continuous functions.

		\item Compute simple limits through algebraic manipulations.

		\item State the Extreme Value Theorem and the Intermediate Value Theorem,
			interpret them geometrically, and use them in simple applications or proofs.
	\end{itemize}

	Notice that the proofs of EVT and IVT are definitely too difficult for this course.

	%%%
	\subsection{3. Derivatives}

	\begin{itemize}
		\item Interpret derivative geometrically (as slope) and physically (as rate
			of change).

		\item Understand and use the formal definition of derivative.

		\item Compute derivatives from the definition.

		\item Compute derivatives of common functions quickly with differentiation rules,
			including implicit differentiation.

		\item Prove the differentiation rules from the definition, and create other
			simple proofs using the definition of derivative. A full proof of the Chain
			Rule is probably too much for this course.

		\item Understand how a function may fail to be differentiable, and identify
			examples geometrically and algebraically.
	\end{itemize}

	%%%
	\subsection{4. Transcendental functions}

	\begin{itemize}
		\item Understand the definitions of function, domain, codomain, range, one-to-one,
			onto, invertible function, and inverse function; compute them in examples,
			and use them in very simple proofs.

		\item Understand the relation between one-to-one and invertible (using the
			common convention in calculus where the codomain is redefined to be the
			range when needed and we never worry about surjectivity).

		\item For non one-to-one functions, work with restrictions to different
			subdomains, construct the corresponding inverses, and understand what the
			compositions will be.

		\item Obtain an identity for the inverse of a function given the derivative of
			the function.

		\item Compute complex derivative using exponentials and logarithms, and use
			``logarithmic differentiation".

		\item Define the inverse trigonometric functions rigorously, compute with
			them, and derive equations for their derivatives.
	\end{itemize}

	%%%
	\subsection{5. The Mean Value Theorem and applications}

	\begin{itemize}
		\item Write definitions and state theorems precisely. This includes the Mean
			Value Theorem, Rolle's Theorem, and theorems about monotonocity, extrema,
			and local extrema. Interpret them geometrically.

		\item Understand and create simple proofs using these theorems.

		\item Find extrema of functions and where they are increasing/decreasing.

		\item Use the MVT and related theorems in simple proofs, calculations, and applications.
	\end{itemize}

	%%%
	\subsection{6. Applications of derivatives and limits}

	\begin{itemize}
		\item Model physical situations using derivatives, and solve ``related-rates
			problems".

		\item Solve applied-optimization problems (i.e.\ translate a real situation into
			an optimization question, an use calculus to optimize it).

		\item When computing limits, determine (and prove) which quotients, products,
			powers, and sums are or are not an indetermine form.

		\item Understand the statement of L'H\^{o}pital's Rule, including all the subtle
			hypotheses that determine when it can or cannot be used.

		\item Compute limits that are in indetermine form, using L'H\^{o}pital's Rule
			and other methods.

		\item State the definition of concavity and interpret it geometrically.

		\item Study the concavity and inflection points of a function.

		\item Define what an asymptote is, interpret it geometrically, and identify them
			in examples.

		\item Sketch the graph of a function, analyzing all the important
			characteristics and points.
	\end{itemize}

	%%%
	\subsection{7. Definition of the integral}

	\begin{itemize}
		\item Use ``sigma notation" for sums.

		\item State the definitions of upper and lower bound, supremum and infinimum
			of a set or a function, and understand equivalent variations of the definitions.
			Compute them in examples, and use them in simple proofs.

		\item Understand all the concepts involve in the formal definition of definite
			integral (partition, upper and lower sum, upper and lower integral, integrable
			function, definite integral), both formally and geometrically. Compute them
			in examples, determine which statements are true or false, and use these concepts
			in simple proofs.

		\item Compute some integrals as limits using Riemann sums.
	\end{itemize}

	%%%
	\subsection{8. The Fundamental Theorem of Calculus}

	\begin{itemize}
		\item Understand the difference between the notions of definite integral,
			antiderivatives, and a function defined as an integral, using the correct
			notation in each case. This includes understanding which notation is correct,
			incorrect (means something else), or nonsense (does not mean anything).
			Know when to use each of these concepts.

		\item Compute elementary antiderivatives.

		\item State the Fundamental Theorem of Calculus (Parts 1 and 2), and use it
			in applications, calculations, and simple proofs.

		\item Compute areas of regions in the plane.
	\end{itemize}

	%%%
	\subsection{9. Integration methods}

	There are two integration methods that are important by themselves: substitution
	and parts. Students need to become fluent in performing these computations, as
	well as understand why the algorithm works, and be able to justify it using
	differentiation rules.

	{\parskip=0.4\baselineskip

	After that, drilling complicated integration methods into students is obsolete. Almost nobody performs them by hand anymore. Instead, when we study other integration methods, we are teaching problem solving, not integration methods. We want students to learn the more general idea of ``reduce this problem to another one I have already solved", which is the essence of all these methods. So we focus on students understanding why the methods work and how someone could come up with them in the first place, rather than applying it to super-complex examples or memorizing a long list of algorithms. To a certain extent it does not matter which other integration methods (if any) we teach. }

	%%%
	\subsection{10. Applications of the integral}

	Many calculus courses study a long list of applications of the integral (volumes,
	surface area, arc length, work, centre of mass, ...) These applications are reduced
	to memorizing a dozen formulas and nothing else. This is of dubious value. Even
	if we restrict ourselves to common and useful applications, if all we do is
	give students a formula to memorize and later ask them to plug different
	functions into it, what is the point? Any time they need one such formula in
	the future they can look it up, and there is no skill (or difficulty) in plugging
	a function into it.


	Instead, our emphasis is on \emph{understanding} why certain quantities can be computed by integrals, and to be able to derive these formulas, including new ones, themselves. That way, whichever application they may need in the future, they will be ready for it. The ideal learning outcome as follows: 

	\begin{quotation}
		``A student takes a second-year class in electricity and magnetism. They
		learn that a charged particle creates an electric potential at all points in
		space. They learn that the electric potential created by various charged particles
		is the sum of the potentials created by each one individually. Based only on
		this (and on what they learned about integrals as limits of Riemann sums in our
		course) they immediately understand how to write the electric potential
		created by a continuous distribution of charges as an integral."
	\end{quotation}

	If we get them to that point, it does not matter which specific applications
	they have ``learned".

	%%%
	\subsection{11. Sequences}

	\begin{itemize}
		\item Understand the (various equivalent) definitions of sequence,
			convergence, various types of divergence, monotonicity, and boundness. We
			do not introduce subsequences.

		\item Transfer all the results about limits of functions to limits of sequences,
			when appropriate.

		\item Use these in computations, examples, and simple proofs.

		\item In particular, know (and be able to justify) the relations between
			them; most notably 1) convergent implies bounded, and 2) bounded $+$
			monotonic implies convergent.

		\item Understand the definition of ``$a_{n}$ grows much slower than $b_{n}$".
			Memorize (and be able to prove) the hierarchy of divergent sequences (logarithms,
			polynomials, exponentials, factorials, and power-exponentials), and use it
			in simple proofs and limit computations.
	\end{itemize}

	%%%
	\subsection{12. Improper integrals}

	\begin{itemize}
		\item Understand that improper integrals are a new concept that we have to define.
			This means that we cannot just assume that improper integrals are always
			well-defined and have the same properties as regular integrals.

		\item State and understand the definition of improper integral and use it to
			compute simple improper integrals.

		\item Memorize the standard family of convergent/divergent improper integrals:
			${\displaystyle \int^{\infty}_1 \frac{dx}{x^{p}}}$.

		\item Understand the statement of the two Comparison Tests (Basic Comparison
			Test and Limit-Comparison Test), why they are true, and use them to prove certain
			improper integrals are convergent and divergent, or in simple proofs.
	\end{itemize}

	%%%
	\subsection{13. Series}

	\begin{itemize}
		\item Never confuse a sequence with a series. Use proper notation for both
			concepts.

		\item Like for improper integrals, understand that an infinite sum is a new
			concept that we have to define.

		\item Understand the definition, be able to explain it, and use it to compute
			the value of simple series (mostly geometric and telescopic series).

		\item Estimate a series numerically.

		\item Know that we cannot take for granted that properties of finite sums carry
			to infinite sums. Know which properties carry and prove so from the
			definition.

		\item State the standard convergence tests and justify intuitively whey they
			work (even if not formally prove them). Correctly use them to determine
			whether a series is convergent (and even conditionally convergent or absolutely
			convergent) or divergent, as well as in simple proofs.
	\end{itemize}

	%%%
	\subsection{14. Power series and Taylor series}

	\begin{itemize}
		\item Identify power series and compute their interval of convergence. %Recognize them as a particular case of a function defined as a series depending on a parameter.

		\item Understand (and use) which properties of polynomials apply to power series
			in the interior of the interval of convergence.

		\item Understand various equivalent definition of Taylor polynomials and Taylor
			series (as a an approximation of a function near a point with an error
			that approaches 0 ``quickly", as a polynomial with the right derivatives, or
			via the standard formula). Compute with them.

		\item Memorize the standard Maclaurin series, and be able to re-derive them.

		\item Use all the common tricks to obtain new Taylor series expansions from old
			ones.

		\item Use Taylor series in many common applications (possibly limits, integrals,
			estimations, adding series, physics, asymptotic behaviour...)
	\end{itemize}

	Throughout this last unit I recommend focusing on conceptual understanding,
	computations, and applications, while avoiding proofs, particularly the most technical
	ones. Understanding, writing, and critiquing simple proofs is one objective
	for this course that has been practiced in all the previous units; however,
	the proofs in this unit are often too technical, hard, or long, so they are not
	appropriate for that learning objective. Most of the aspects of the theory and
	most of the theorems are illustrated by how we use them in the applications.
\end{document}

%============
%============
%============
%============
